% BEGIN REV09_FORMA_PARTIDA_ENERGIA
\section{Split form, integrality, and lattice energy}
\label{corpus9:xii:forma-partida}

The deformation in Theorem~\ref{rec:thm:red} starts from the marked
lattice \(\Lambda\), the direction \(u=P_3P_{11}K\), and the
previously generated regional coordinates. Let \(E\) be its
twenty-four-dimensional real space, and let
\(\langle\,\cdot\,,\,\cdot\,\rangle_B\) denote the positive metric
of the twelve orthogonal planes of type \(A_2\). We retain the
normalization of \eqref{rec:eq:flujo} and \eqref{rec:eq:red}:
\begin{equation}
 \kappa=\frac{\log\rho}{\|u\|},\qquad h_i=\kappa u_i,
 \qquad
 G_t=\bigoplus_{i=1}^{12}e^{th_i}I_{A_{2,i}},\qquad
 \Lambda_t=G_t\Lambda.
 \label{corpus9:xii:normalizacion-flujo}
\end{equation}
The operator \(G_t\) is self-adjoint for \(B\), and
\(G_t^{-1}=G_{-t}\). If the same collection is written with the
unnormalized parameter, \(D_s=\bigoplus_i e^{su_i}I_{A_{2,i}}\), the
exact relation is \(G_t=D_{\kappa t}\). The deformation parameter and
the TPK clock retain their respective domains.

\subsection{Arithmetic variation of the Euclidean projection}
\label{corpus9:xii:variacion-euclidea}

The ordered chart in Lemma~\ref{lem:k-sector-no-nulo} determines,
through the finite sum \eqref{eq:k-proyectores-incidencia},
\begin{equation}
 u_1=-\frac{495}{4}-\frac{233}{10}\sqrt5,
 \qquad \sum_{i=1}^{12}u_i=0.
 \label{corpus9:xii:coordenada-testigo}
\end{equation}
The regional normalization satisfies
\(\rho=\pi_{\HMT}/\varphi_{\HMT}^2>1\): in the realization already
established, \(\pi_{\HMT}>3\) and
\(\varphi_{\HMT}^2=(3+\sqrt5)/2<3\). Thus \(\kappa>0\).
These comparisons concern the regional outputs after their generation.

\begin{proposition}[Local witness for variation of integrality]
\label{corpus9:xii:testigo-integralidad}
There exists \(\delta>0\) such that, for \(0<|t|<\delta\), the
lattice \(\Lambda_t\) contains a vector with noninteger squared norm.
Consequently, its Euclidean metric lacks integrality and an isometry
with the Leech lattice.
\end{proposition}
\begin{proof}
The marked vector in \eqref{exc:vector} is
\(v_\alpha=4\rho_1+\sum_{j=2}^{12}\rho_j\), with
\(\|\rho_j\|_B^2=2\) in mutually orthogonal planes and
\(v_\alpha^2=54\). The neighbour construction
\eqref{exc:vecino} gives \(v_\alpha/3\in\Lambda\). Its transported
norm is
\begin{equation}
 f(t)=\left\|G_t\frac{v_\alpha}{3}\right\|_B^2
 =\frac29\left(16e^{2th_1}+\sum_{j=2}^{12}e^{2th_j}\right).
 \label{corpus9:xii:norma-testigo}
\end{equation}
The zero sum and \eqref{corpus9:xii:coordenada-testigo} give
\begin{equation}
 \begin{aligned}
 f(0)&=6,\\
 f'(0)&=\frac49\left(16h_1+\sum_{j=2}^{12}h_j\right)
       =\frac{20}{3}\kappa u_1
       =\kappa\left(-825-\frac{466}{3}\sqrt5\right)\ne0.
 \end{aligned}
 \label{corpus9:xii:derivada-testigo}
\end{equation}
Continuity of \(f'\) provides an interval
\((-\delta,\delta)\) on which \(f\) is strictly monotone.
Reducing \(\delta\) also gives \(11/2<f(t)<13/2\). For
\(0<|t|<\delta\), monotonicity yields \(f(t)\ne f(0)=6\);
six is the only integer in that interval. Thus \(G_t(v_\alpha/3)\)
has noninteger squared norm. An integral lattice has integer squared
norms, and isometries preserve those norms. Integrality of the Leech
lattice proves the final conclusion.
\end{proof}

The deformation preserves covolume and ternary symmetry while varying
its arithmetic metric structure. Reciprocity
\(\Lambda_t^*=\Lambda_{-t}\) uses self-adjointness of the flow and
initial self-duality. An ambient permutation interchanging expansions
and contractions acts on the glued lattice when it also preserves the
code and the neighbour from which that lattice is constructed.

\subsection{Integral self-duality for the split-signature form}
\label{corpus9:xii:autodualidad-partida}

Simultaneous consideration of both projections uses the space
\(E\oplus E\), with bilinear form
\begin{equation}
 \mathcal B((x,p),(x',p'))
 =\langle x,p'\rangle_B+\langle p,x'\rangle_B.
 \label{corpus9:xii:forma-bilineal-partida}
\end{equation}
Define
\begin{equation}
 \begin{aligned}
 \mathcal G_t&=G_t\oplus G_{-t},\\
 \Gamma_t&=\mathcal G_t(\Lambda\oplus\Lambda)
 =\{(G_t\lambda,G_{-t}\mu):\lambda,\mu\in\Lambda\}.
 \end{aligned}
 \label{corpus9:xii:red-partida}
\end{equation}

\begin{theorem}[Integrality, evenness, and self-duality of the split lattice]
\label{corpus9:xii:teorema-partido}
The form \(\mathcal B\) has signature \((24,24)\).
For every \(t\in\mathbb R\), \(\mathcal G_t\) is an isometry of
\(\mathcal B\), and \(\Gamma_t\) is an integral, even, self-dual
lattice with respect to that form.
\end{theorem}
\begin{proof}
The coordinates
\[
 p_L=\frac{x+p}{\sqrt2},\qquad p_R=\frac{x-p}{\sqrt2}
\]
give an invertible linear change, with
\(x=(p_L+p_R)/\sqrt2\) and \(p=(p_L-p_R)/\sqrt2\). They satisfy
\[
 \mathcal B((x,p),(x,p))=\|p_L\|_B^2-\|p_R\|_B^2.
\]
Positivity of \(B\) in dimension twenty-four gives signature
\((24,24)\). Self-adjointness of \(G_t\) and the identity
\(G_tG_{-t}=I\) imply
\[
 \langle G_tx,G_{-t}p'\rangle_B
 =\langle x,G_tG_{-t}p'\rangle_B=\langle x,p'\rangle_B.
\]
The other cross product satisfies the same equality. Hence
\(\mathcal B(\mathcal G_t\xi,\mathcal G_t\eta)
=\mathcal B(\xi,\eta)\).

On \(\Gamma_0=\Lambda\oplus\Lambda\), the cross products are
integers because \(\Lambda\) is integral, and
\[
 \mathcal B((\lambda,\mu),(\lambda,\mu))
 =2\langle\lambda,\mu\rangle_B\in2\mathbb Z.
\]
This proves integrality and evenness. If \((x,p)\) belongs to the
dual of \(\Gamma_0\) with respect to \(\mathcal B\), its pairings
with \((\lambda,0)\) and \((0,\mu)\), for all
\(\lambda,\mu\in\Lambda\), require respectively
\(p,x\in\Lambda^*=\Lambda\). Conversely, these two memberships
ensure that every pairing is integral. Thus
\(\Gamma_0^{*\mathcal B}=\Gamma_0\). The isometry
\(\mathcal G_t\) transports all three properties, giving
\(\Gamma_t^{*\mathcal B}=\mathcal G_t\Gamma_0^{*\mathcal B}
=\Gamma_t\).
\end{proof}

The metric determines the domain of each conclusion. The cross products
of \(\mathcal B\) preserve integrality of \(\Gamma_t\), including
at the parameter values of
Proposition~\ref{corpus9:xii:testigo-integralidad}. The lattices
\(\Gamma_t\), equipped with that split form, are mutually isometric.
Their Euclidean projections \(\Lambda_t\) and \(\Lambda_{-t}\),
equipped with \(B\), vary with the parameter.

\begin{proposition}[Exchange of reciprocal projections]
\label{corpus9:xii:involucion-partida}
The involution \(\mathcal J(x,p)=(p,x)\) satisfies
\[
 \mathcal J\Gamma_t=\Gamma_{-t},\qquad
 \mathcal J\mathcal G_t\mathcal J=\mathcal G_{-t}.
\]
In the coordinates \((p_L,p_R)\), it fixes \(p_L\) and reverses
the sign of \(p_R\).
\end{proposition}
\begin{proof}
The exchange transforms \((G_t\lambda,G_{-t}\mu)\) into
\((G_{-t}\mu,G_t\lambda)\). As \(\lambda,\mu\) range over
\(\Lambda\), the image ranges over \(\Gamma_{-t}\). Conjugation
interchanges the two blocks of \(\mathcal G_t\), producing
\(\mathcal G_{-t}\). Finally, \(x+p\) is preserved and \(x-p\)
changes sign.
\end{proof}

This involution acts on the complete lattice construction. A subsequent
physical representation must also transport the operators and
identification data proper to that realization.

\subsection{Positive operator and unitary equivalence of domains}
\label{corpus9:xii:energia-reticular}

The energy of the two projections is defined on the fixed lattice by
\begin{equation}
 \mathsf E_t^{\rm lat}(\lambda,\mu)
 =\|G_t\lambda\|_B^2+\|G_{-t}\mu\|_B^2,
 \qquad \lambda,\mu\in\Lambda.
 \label{corpus9:xii:energia}
\end{equation}
It is nonnegative, vanishes exactly at \((0,0)\), and satisfies
\(\mathsf E_t^{\rm lat}(\lambda,\mu)
=\mathsf E_{-t}^{\rm lat}(\mu,\lambda)\). Exchange of the two
contributions expresses, on the twelve marked planes, the reciprocity
of a sum of oppositely scaled terms. For \(M_h=\max_i|h_i|\), the
orthogonal decomposition into planes gives
\begin{equation}
 e^{-2|t|M_h}\|v\|_B^2\leq\|G_tv\|_B^2
                    \leq e^{2|t|M_h}\|v\|_B^2.
 \label{corpus9:xii:cotas-cuadraticas}
\end{equation}
Indeed, each summand \(\|v_i\|_B^2\) is multiplied by
\(e^{2th_i}\), which lies between these bounds.

\begin{theorem}[Self-adjointness, compact resolvent, and unitary duality]
\label{corpus9:xii:operador-positivo}
On \(\ell^2(\Lambda\oplus\Lambda)\), the operator
\[
 (H_t^{\rm lat}\psi)(\lambda,\mu)
 =\mathsf E_t^{\rm lat}(\lambda,\mu)\psi(\lambda,\mu)
\]
with maximal domain
\begin{equation}
 \mathcal D(H_t^{\rm lat})=
 \left\{\psi\in\ell^2(\Lambda\oplus\Lambda):
 \sum_{\lambda,\mu}\bigl(\mathsf E_t^{\rm lat}(\lambda,\mu)\bigr)^2
                  |\psi(\lambda,\mu)|^2<\infty\right\}
 \label{corpus9:xii:dominio-operador}
\end{equation}
is positive and self-adjoint with compact resolvent. The map
\(U\psi(\lambda,\mu)=\psi(\mu,\lambda)\) is unitary and satisfies
\begin{equation}
 U\mathcal D(H_t^{\rm lat})=\mathcal D(H_{-t}^{\rm lat}),\qquad
 UH_t^{\rm lat}U^{-1}=H_{-t}^{\rm lat}.
 \label{corpus9:xii:equivalencia-unitaria}
\end{equation}
\end{theorem}
\begin{proof}
Finitely supported functions belong to the domain and form a dense
subspace. The energy is real. If \(\psi\) belongs to the adjoint
domain, evaluating its functional against each singleton-supported
orthonormal basis vector forces the corresponding adjoint coordinate
to be \(\mathsf E_t^{\rm lat}(\lambda,\mu)\psi(\lambda,\mu)\).
Membership of this sequence in \(\ell^2\) is exactly condition
\eqref{corpus9:xii:dominio-operador}. Conversely, that condition permits
pairing with every function in the domain by Cauchy--Schwarz. The
adjoint domain and action coincide with those of \(H_t^{\rm lat}\),
proving self-adjointness. Moreover,
\[
 \langle\psi,H_t^{\rm lat}\psi\rangle
 =\sum_{\lambda,\mu}\mathsf E_t^{\rm lat}(\lambda,\mu)
                         |\psi(\lambda,\mu)|^2\geq0.
\]
By \eqref{corpus9:xii:cotas-cuadraticas},
\[
 \mathsf E_t^{\rm lat}(\lambda,\mu)
 \geq e^{-2|t|M_h}
       \bigl(\|\lambda\|_B^2+\|\mu\|_B^2\bigr).
\]
Discreteness of \(\Lambda\oplus\Lambda\) implies that every energy
sublevel set is finite. The resolvent \((H_t^{\rm lat}+I)^{-1}\)
is diagonal, with coefficients \((1+\mathsf E_t^{\rm lat})^{-1}\).
Its restriction to the sublevel \(\mathsf E_t^{\rm lat}\leq R\)
has finite rank, and the remainder has norm at most \((1+R)^{-1}\).
The resolvent is therefore a norm limit of finite-rank operators and
is compact.

Index exchange preserves the \(\ell^2\) norm and satisfies
\(U^{-1}=U\). Replacing \((\lambda,\mu)\) by \((\mu,\lambda)\)
in the sum defining the domain and using the energy equality gives
the first identity in \eqref{corpus9:xii:equivalencia-unitaria}.
On that domain,
\[
 (UH_t^{\rm lat}U^{-1}\psi)(\lambda,\mu)
 =\mathsf E_t^{\rm lat}(\mu,\lambda)\psi(\lambda,\mu)
 =\mathsf E_{-t}^{\rm lat}(\lambda,\mu)\psi(\lambda,\mu),
\]
which proves the operator identity together with its domains.
\end{proof}

The three preservation laws are thus typed. The Euclidean metric
preserves covolume and ternary symmetry and relates opposite parameters
through the dual lattice. The split form preserves integrality,
evenness, and self-duality of the doubled lattice. The operator
realization preserves energy through a unitary equivalence that
transports its domain. The theta reader in
Theorem~\ref{rec:thm:theta} expresses reciprocity in the convergent sum
of lengths, and Mellin continuation preserves its polar structure by
Theorem~\ref{rec:thm:mellin}.
% END REV09_FORMA_PARTIDA_ENERGIA
